% PRISM real-world experiments section draft for a CoRL-style paper.
% Required packages in the main paper preamble:
% \usepackage{booktabs}

\section{Real-World Experiments}
\label{sec:real_world}

We evaluate PRISM on a wheeled dual-arm mobile manipulator equipped with one overhead RGB camera and two wrist-mounted RGB cameras. The policy input consists of synchronized multi-view images together with the robot's proprioceptive state, including the base, both arms, and grippers as available on the platform. The state-prefix signature is computed online from the executed state trajectory using the same update rule as in Section~\ref{sec:method_inference}, so the controller never observes future states. Unless otherwise stated, we instantiate PRISM on top of the ACT action expert and compare against baselines that isolate the contribution of state-prefix conditioning, fixed-budget visual prefix memory, and signature-indexed routing.

Our real-world evaluation is designed to test delayed disambiguation under long-horizon partial observability rather than instantaneous visual recognition. To make the tasks faithful to the claim of PRISM, we impose three design principles. First, the task-relevant cue is injected early and is short-lived: it is available only during an inspection or preload phase. Second, the episode then enters a canonical merge stage in which different branches are brought to nearly identical current observations and robot poses. Third, the final decision is irreversible, so success depends on retaining the executed prefix rather than reacting to the last frame. We additionally sweep the length of the shared post-merge gap to study how performance degrades with memory horizon.

\begin{table}[t]
    \centering
    \small
    \caption{Summary of the real-world PRISM tasks. Each task contains an early cue, a shared merge stage, and a delayed irreversible decision.}
    \label{tab:real_world_tasks}
    \begin{tabular}{p{2.8cm}p{3.0cm}p{3.1cm}p{3.0cm}}
        \toprule
        Task & Early cue & Delayed decision & Primary capability \\
        \midrule
        Desktop cleaning & brief inspection of a hidden contamination type & select and apply the correct tool after occlusion and return to a canonical pose & delayed visual disambiguation \\
        Symmetric port insertion & four-state bimanual preload encoded in the robot state & choose the correct symmetric port after equalization & state-prefix routing and delayed branch memory \\
        Sequential dual-query servicing & ordered inspection of two hidden service sites & recover the correct two tools and service the sites in order & multi-event and order-sensitive memory \\
        Rescaled pattern reproduction & brief observation of a reference unlock trajectory & reproduce the node sequence on a grid of different size & geometric transfer under task-frame rescaling \\
        \bottomrule
    \end{tabular}
\end{table}

\subsection{Platform and Task Design}
\label{sec:real_world_setup}

All physical tasks are executed with the same sensing stack, policy interface, and online memory update. This shared setup ensures that any performance gap is not attributable to task-specific engineering. We randomize object placement, camera-visible clutter, tool arrangement, and post-merge delay within a predefined range for every episode. We also verify the intended ambiguity of each task before policy training: at the delayed decision stage, a lightweight probe classifier that only sees the current observation should remain close to chance in predicting the hidden branch identity. This sanity check rules out trivial shortcuts from the last frame and makes the real-world evaluation better aligned with the motivation of PRISM.

\subsection{Task Suite}
\label{sec:real_world_tasks}

\paragraph{Delayed Tool Selection for Desktop Cleaning.}
This task refines the initial ``desktop cleaning'' setup into a delayed tool-selection problem. Several visually similar work zones are placed on the desk, and one zone contains a contamination type sampled from a small tool dictionary, such as a liquid spill, dry particles, or an erasable marker trace. The robot is first allowed to inspect the relevant zone with the overhead or wrist cameras. After inspection, the contaminated region is hidden by an opaque cover, and the robot returns to a canonical intermediate pose that matches across all contamination types. From this shared pose, the robot must move to a tool rack, choose the correct cleaning tool, and then return to clean the hidden target region. The key point is that tool selection occurs after the contamination is no longer directly visible, so the task cannot be solved by instantaneous image classification alone. We report tool-selection accuracy, task completion rate, residual contamination, and success as a function of the post-merge gap length.

\paragraph{Four-State Preload to Symmetric Port Insertion.}
This task is the most direct real-world probe of state-prefix memory. The two end-effectors hold a semi-rigid probe that is preloaded into one of four initial bimanual states defined by left/right high-low start slots. These four reset states produce distinct early joint and end-effector trajectories, hence distinct state-prefix signatures, but they are subsequently funneled through an equalization guide that brings the robot to a common intermediate pose with a nearly identical visual appearance. From that merged state, the robot must choose one of four symmetric target ports and complete insertion, where the correct port is determined only by the initial preload code. This task is therefore explicitly constructed so that the delayed branching decision should depend on the executed state prefix rather than on the current image. We report final port classification accuracy, insertion success rate, and, when the fixture exposes explicit intermediate branch gates, gate-wise accuracy before the terminal insertion.

\paragraph{Sequential Dual-Query Servicing.}
To test whether PRISM can retain more than a single early cue, we introduce a two-query sequential servicing task. At the beginning of the episode, the robot briefly inspects two service sites in sequence, where each site corresponds to a different hidden servicing requirement and therefore a different tool. The sites are then covered, and the robot executes a shared transfer segment that equalizes the current pose before any tool is selected. The robot must thereafter retrieve the two correct tools and service the two sites in the correct temporal order. This setting is deliberately harder than single-query desktop cleaning because the policy must remember both content and order, not just a single latent task label. We report first-query accuracy, second-query accuracy, order accuracy, and full-episode success. Relative to the single-step tasks, this experiment is particularly useful for distinguishing a structured slot-based memory from baselines that compress history into a single hidden vector.

\paragraph{Rescaled Pattern Reproduction.}
Our original ``unlock pattern'' idea is retained as a supplementary geometry transfer task, but we frame it as generalization under task-frame rescaling rather than as evidence of an intrinsic scale-invariance property. The robot first observes a short reference trajectory on a $3 \times 3$ point grid and then, after a shared motion segment, must reproduce the same node sequence on a target grid whose spatial scale differs from the reference. To avoid confounds from trivial pixel-level resizing, we express the target in a task-centered coordinate frame and evaluate both seen and unseen grid scales. This task does not replace the memory-centric tasks above; instead, it tests whether a prefix-conditioned controller can preserve and replay a geometric instruction when the execution workspace is rescaled. We report exact node-sequence accuracy, normalized trajectory error, illegal-transition rate, and success on held-out scales.

\subsection{Baselines, Metrics, and Protocol}
\label{sec:real_world_protocol}

Unless otherwise stated, each real-world task compares the following ACT-based controllers: (i) Vanilla ACT without any explicit history module; (ii) ACT+Signature, which receives the current state-prefix signature token but no streaming visual prefix memory; (iii) ACT+Prefix Memory, which uses the same fixed-budget prefix memory as PRISM but removes signature-indexed routing; (iv) ACT+First-Frame Anchor, a strong short-horizon baseline that preserves the initial visual cue without modeling the full prefix; and (v) PRISM. This baseline suite separates three questions: whether the task truly requires history, whether state-prefix information alone is sufficient, and whether the main benefit of PRISM comes from memory capacity or from signature-indexed memory routing.

For every task, we report overall success together with stage-wise metrics that localize where the failure occurs. These stage-wise metrics include tool selection accuracy for desktop cleaning, port selection accuracy for symmetric insertion, first/second-query and order accuracy for dual-query servicing, and sequence-level geometric metrics for pattern reproduction. We additionally evaluate success as a function of the post-merge gap length, for example by sweeping delays such as $\{0, 3, 6, 10\}$~s or an equivalent shared-motion budget. This memory-gap sweep is important because it reveals whether a method merely exploits short-horizon cues or genuinely preserves decision-relevant prefix information over extended execution.

Finally, we treat the real-world suite not only as a benchmark of end-task success but also as a controlled causal test of the PRISM design. If Vanilla ACT already solves the delayed decision stage, then the task is insufficiently ambiguous. If ACT+Prefix Memory matches PRISM, then signature-indexed routing is unnecessary for that setting. Conversely, consistent gains of PRISM over both ACT+Signature and ACT+Prefix Memory would support the central claim of the paper: preserving the executed prefix is not enough by itself, and structuring the memory update by the state-prefix signature yields a more reliable long-horizon controller in the real world.

% TODO:
% 1. Insert the number of demonstrations and evaluation trials per task.
% 2. State whether the mobile base is fixed after reset or actively controlled.
% 3. Add concrete tool categories, contamination definitions, and acceptance thresholds.
% 4. Add the exact train/test split over layouts, delays, and object identities.
